% PYNE Paper 03 — Series History Management and Runtime Performance
% arXiv target: cs.DS / cs.PF
\documentclass[11pt,a4paper]{article}

\usepackage[margin=1in]{geometry}
\usepackage[T1]{fontenc}
\usepackage[utf8]{inputenc}

% Shared macros for PYNE arXiv paper series
% Include with: % Shared macros for PYNE arXiv paper series
% Include with: % Shared macros for PYNE arXiv paper series
% Include with: \input{../common/macros.tex}

\usepackage{amsmath,amssymb,amsfonts}
\usepackage{amsthm}
\usepackage{booktabs}
\usepackage{graphicx}
\usepackage{xcolor}
\usepackage{hyperref}
\usepackage{url}
\usepackage{algorithm}
\usepackage{algpseudocode}
\usepackage{listings}
\usepackage{tikz}
\usepackage{pgfplots}
\pgfplotsset{compat=1.17}
\usetikzlibrary{arrows.meta,positioning,fit,calc,shapes.geometric,shapes.multipart,decorations.pathreplacing,backgrounds,chains,matrix,patterns}
\usepackage{subcaption}
\usepackage{multirow}
\usepackage{tabularx}
\usepackage{enumitem}
\usepackage{microtype}
\usepackage{balance}
\usepackage{csquotes}
\usepackage{natbib}

% Colors
\definecolor{pyneorange}{HTML}{F97316}
\definecolor{pyneblue}{HTML}{2563EB}
\definecolor{pynegray}{HTML}{6B7280}
\definecolor{pynegreen}{HTML}{16A34A}
\definecolor{codebg}{HTML}{F8FAFC}
\definecolor{codekw}{HTML}{7C3AED}
\definecolor{codestr}{HTML}{B45309}
\definecolor{codecomment}{HTML}{64748B}

\hypersetup{
  colorlinks=true,
  linkcolor=pyneblue,
  citecolor=pynegreen,
  urlcolor=pyneorange,
  pdftitle={},
  pdfauthor={HOOX / PYNE Project}
}

% Code listings
\lstdefinestyle{pine}{
  basicstyle=\ttfamily\footnotesize,
  backgroundcolor=\color{codebg},
  keywordstyle=\color{codekw}\bfseries,
  stringstyle=\color{codestr},
  commentstyle=\color{codecomment}\itshape,
  breaklines=true,
  frame=single,
  rulecolor=\color{black!20},
  numbers=left,
  numberstyle=\tiny\color{pynegray},
  numbersep=6pt,
  xleftmargin=1.5em,
  framexleftmargin=1.2em,
  showstringspaces=false,
  tabsize=2,
  morekeywords={indicator,strategy,library,plot,plotshape,fill,hline,var,varip,if,else,for,while,switch,true,false,na,and,or,not,import,export,type,method,enum},
  morecomment=[l]{//},
  morestring=[b]",
}
\lstdefinestyle{python}{
  language=Python,
  basicstyle=\ttfamily\footnotesize,
  backgroundcolor=\color{codebg},
  keywordstyle=\color{codekw}\bfseries,
  stringstyle=\color{codestr},
  commentstyle=\color{codecomment}\itshape,
  breaklines=true,
  frame=single,
  rulecolor=\color{black!20},
  numbers=left,
  numberstyle=\tiny\color{pynegray},
  numbersep=6pt,
  xleftmargin=1.5em,
  framexleftmargin=1.2em,
  showstringspaces=false,
  tabsize=4,
}
\lstset{style=python}

% Theorem environments
\theoremstyle{plain}
\newtheorem{theorem}{Theorem}[section]
\newtheorem{lemma}[theorem]{Lemma}
\newtheorem{proposition}[theorem]{Proposition}
\newtheorem{corollary}[theorem]{Corollary}
\theoremstyle{definition}
\newtheorem{definition}[theorem]{Definition}
\newtheorem{example}[theorem]{Example}
\newtheorem{remark}[theorem]{Remark}

% Shortcuts
\newcommand{\pyne}{\textsc{Pyne}}
\newcommand{\pine}{Pine Script\texttrademark}
\newcommand{\tv}{TradingView\textregistered}
\newcommand{\asdl}{\textsc{Asdl}}
\newcommand{\antlr}{\textsc{Antlr}4}
\newcommand{\numba}{\textsc{Numba}}
\newcommand{\lsp}{\textsc{Lsp}}
\newcommand{\ohlcv}{\textsc{Ohlcv}}
\newcommand{\udt}{\textsc{Udt}}
\newcommand{\jit}{\textsc{Jit}}
\newcommand{\api}{\textsc{Api}}
\newcommand{\cli}{\textsc{Cli}}
% AST acronym (text mode). Keep amsmath's math asterisk \ast intact.
\DeclareRobustCommand{\Ast}{\textsc{Ast}}
% Papers in this series write \ast{} for the AST acronym in prose.
% Redefine robustly; use \mathast for a binary asterisk in formulas.
\let\mathast\ast
\DeclareRobustCommand{\ast}{\textsc{Ast}}
\newcommand{\ir}{\textsc{Ir}}
\newcommand{\ta}{\texttt{ta}}
\newcommand{\code}[1]{\texttt{#1}}
\newcommand{\file}[1]{\texttt{#1}}
\newcommand{\eg}{e.g.}
\newcommand{\ie}{i.e.}
\newcommand{\etal}{\textit{et~al.}}
\newcommand{\resp}{respectively}

% Math helpers
\newcommand{\R}{\mathbb{R}}
\newcommand{\N}{\mathbb{N}}
\newcommand{\Z}{\mathbb{Z}}
% Text-mode NA token for prose; use $\namath$ inside formulas if needed.
\DeclareRobustCommand{\na}{\textsf{na}}
\newcommand{\namath}{\mathsf{na}}
\newcommand{\series}[1]{\mathbf{#1}}
\newcommand{\baridx}{i}
\newcommand{\bars}{n}
\newcommand{\period}{p}
\newcommand{\abserr}{\varepsilon_{\mathrm{abs}}}
\newcommand{\relerr}{\varepsilon_{\mathrm{rel}}}
\newcommand{\maxerr}{\varepsilon_{\mathrm{max}}}
\newcommand{\meanerr}{\varepsilon_{\mathrm{mean}}}

% Pipeline notation — use inside math: $\pipeline$
\newcommand{\pipeline}{\mathrm{src}\!\to\!\mathrm{lex}\!\to\!\mathrm{parse}\!\to\!\mathrm{AST}\!\to\!\mathrm{eval}}

% Disclaimer footnote helper
\newcommand{\trademarknote}{%
  \pine{} and \tv{} are trademarks of TradingView, Inc.
  \pyne{} is an independent, unofficial implementation and is not affiliated with,
  authorized by, sponsored by, or endorsed by TradingView, Inc.
}

% Figure path helper (papers live one level under .papers/)
\graphicspath{{figures/}{../common/figures/}}


\usepackage{amsmath,amssymb,amsfonts}
\usepackage{amsthm}
\usepackage{booktabs}
\usepackage{graphicx}
\usepackage{xcolor}
\usepackage{hyperref}
\usepackage{url}
\usepackage{algorithm}
\usepackage{algpseudocode}
\usepackage{listings}
\usepackage{tikz}
\usepackage{pgfplots}
\pgfplotsset{compat=1.17}
\usetikzlibrary{arrows.meta,positioning,fit,calc,shapes.geometric,shapes.multipart,decorations.pathreplacing,backgrounds,chains,matrix,patterns}
\usepackage{subcaption}
\usepackage{multirow}
\usepackage{tabularx}
\usepackage{enumitem}
\usepackage{microtype}
\usepackage{balance}
\usepackage{csquotes}
\usepackage{natbib}

% Colors
\definecolor{pyneorange}{HTML}{F97316}
\definecolor{pyneblue}{HTML}{2563EB}
\definecolor{pynegray}{HTML}{6B7280}
\definecolor{pynegreen}{HTML}{16A34A}
\definecolor{codebg}{HTML}{F8FAFC}
\definecolor{codekw}{HTML}{7C3AED}
\definecolor{codestr}{HTML}{B45309}
\definecolor{codecomment}{HTML}{64748B}

\hypersetup{
  colorlinks=true,
  linkcolor=pyneblue,
  citecolor=pynegreen,
  urlcolor=pyneorange,
  pdftitle={},
  pdfauthor={HOOX / PYNE Project}
}

% Code listings
\lstdefinestyle{pine}{
  basicstyle=\ttfamily\footnotesize,
  backgroundcolor=\color{codebg},
  keywordstyle=\color{codekw}\bfseries,
  stringstyle=\color{codestr},
  commentstyle=\color{codecomment}\itshape,
  breaklines=true,
  frame=single,
  rulecolor=\color{black!20},
  numbers=left,
  numberstyle=\tiny\color{pynegray},
  numbersep=6pt,
  xleftmargin=1.5em,
  framexleftmargin=1.2em,
  showstringspaces=false,
  tabsize=2,
  morekeywords={indicator,strategy,library,plot,plotshape,fill,hline,var,varip,if,else,for,while,switch,true,false,na,and,or,not,import,export,type,method,enum},
  morecomment=[l]{//},
  morestring=[b]",
}
\lstdefinestyle{python}{
  language=Python,
  basicstyle=\ttfamily\footnotesize,
  backgroundcolor=\color{codebg},
  keywordstyle=\color{codekw}\bfseries,
  stringstyle=\color{codestr},
  commentstyle=\color{codecomment}\itshape,
  breaklines=true,
  frame=single,
  rulecolor=\color{black!20},
  numbers=left,
  numberstyle=\tiny\color{pynegray},
  numbersep=6pt,
  xleftmargin=1.5em,
  framexleftmargin=1.2em,
  showstringspaces=false,
  tabsize=4,
}
\lstset{style=python}

% Theorem environments
\theoremstyle{plain}
\newtheorem{theorem}{Theorem}[section]
\newtheorem{lemma}[theorem]{Lemma}
\newtheorem{proposition}[theorem]{Proposition}
\newtheorem{corollary}[theorem]{Corollary}
\theoremstyle{definition}
\newtheorem{definition}[theorem]{Definition}
\newtheorem{example}[theorem]{Example}
\newtheorem{remark}[theorem]{Remark}

% Shortcuts
\newcommand{\pyne}{\textsc{Pyne}}
\newcommand{\pine}{Pine Script\texttrademark}
\newcommand{\tv}{TradingView\textregistered}
\newcommand{\asdl}{\textsc{Asdl}}
\newcommand{\antlr}{\textsc{Antlr}4}
\newcommand{\numba}{\textsc{Numba}}
\newcommand{\lsp}{\textsc{Lsp}}
\newcommand{\ohlcv}{\textsc{Ohlcv}}
\newcommand{\udt}{\textsc{Udt}}
\newcommand{\jit}{\textsc{Jit}}
\newcommand{\api}{\textsc{Api}}
\newcommand{\cli}{\textsc{Cli}}
% AST acronym (text mode). Keep amsmath's math asterisk \ast intact.
\DeclareRobustCommand{\Ast}{\textsc{Ast}}
% Papers in this series write \ast{} for the AST acronym in prose.
% Redefine robustly; use \mathast for a binary asterisk in formulas.
\let\mathast\ast
\DeclareRobustCommand{\ast}{\textsc{Ast}}
\newcommand{\ir}{\textsc{Ir}}
\newcommand{\ta}{\texttt{ta}}
\newcommand{\code}[1]{\texttt{#1}}
\newcommand{\file}[1]{\texttt{#1}}
\newcommand{\eg}{e.g.}
\newcommand{\ie}{i.e.}
\newcommand{\etal}{\textit{et~al.}}
\newcommand{\resp}{respectively}

% Math helpers
\newcommand{\R}{\mathbb{R}}
\newcommand{\N}{\mathbb{N}}
\newcommand{\Z}{\mathbb{Z}}
% Text-mode NA token for prose; use $\namath$ inside formulas if needed.
\DeclareRobustCommand{\na}{\textsf{na}}
\newcommand{\namath}{\mathsf{na}}
\newcommand{\series}[1]{\mathbf{#1}}
\newcommand{\baridx}{i}
\newcommand{\bars}{n}
\newcommand{\period}{p}
\newcommand{\abserr}{\varepsilon_{\mathrm{abs}}}
\newcommand{\relerr}{\varepsilon_{\mathrm{rel}}}
\newcommand{\maxerr}{\varepsilon_{\mathrm{max}}}
\newcommand{\meanerr}{\varepsilon_{\mathrm{mean}}}

% Pipeline notation — use inside math: $\pipeline$
\newcommand{\pipeline}{\mathrm{src}\!\to\!\mathrm{lex}\!\to\!\mathrm{parse}\!\to\!\mathrm{AST}\!\to\!\mathrm{eval}}

% Disclaimer footnote helper
\newcommand{\trademarknote}{%
  \pine{} and \tv{} are trademarks of TradingView, Inc.
  \pyne{} is an independent, unofficial implementation and is not affiliated with,
  authorized by, sponsored by, or endorsed by TradingView, Inc.
}

% Figure path helper (papers live one level under .papers/)
\graphicspath{{figures/}{../common/figures/}}


\usepackage{amsmath,amssymb,amsfonts}
\usepackage{amsthm}
\usepackage{booktabs}
\usepackage{graphicx}
\usepackage{xcolor}
\usepackage{hyperref}
\usepackage{url}
\usepackage{algorithm}
\usepackage{algpseudocode}
\usepackage{listings}
\usepackage{tikz}
\usepackage{pgfplots}
\pgfplotsset{compat=1.17}
\usetikzlibrary{arrows.meta,positioning,fit,calc,shapes.geometric,shapes.multipart,decorations.pathreplacing,backgrounds,chains,matrix,patterns}
\usepackage{subcaption}
\usepackage{multirow}
\usepackage{tabularx}
\usepackage{enumitem}
\usepackage{microtype}
\usepackage{balance}
\usepackage{csquotes}
\usepackage{natbib}

% Colors
\definecolor{pyneorange}{HTML}{F97316}
\definecolor{pyneblue}{HTML}{2563EB}
\definecolor{pynegray}{HTML}{6B7280}
\definecolor{pynegreen}{HTML}{16A34A}
\definecolor{codebg}{HTML}{F8FAFC}
\definecolor{codekw}{HTML}{7C3AED}
\definecolor{codestr}{HTML}{B45309}
\definecolor{codecomment}{HTML}{64748B}

\hypersetup{
  colorlinks=true,
  linkcolor=pyneblue,
  citecolor=pynegreen,
  urlcolor=pyneorange,
  pdftitle={},
  pdfauthor={HOOX / PYNE Project}
}

% Code listings
\lstdefinestyle{pine}{
  basicstyle=\ttfamily\footnotesize,
  backgroundcolor=\color{codebg},
  keywordstyle=\color{codekw}\bfseries,
  stringstyle=\color{codestr},
  commentstyle=\color{codecomment}\itshape,
  breaklines=true,
  frame=single,
  rulecolor=\color{black!20},
  numbers=left,
  numberstyle=\tiny\color{pynegray},
  numbersep=6pt,
  xleftmargin=1.5em,
  framexleftmargin=1.2em,
  showstringspaces=false,
  tabsize=2,
  morekeywords={indicator,strategy,library,plot,plotshape,fill,hline,var,varip,if,else,for,while,switch,true,false,na,and,or,not,import,export,type,method,enum},
  morecomment=[l]{//},
  morestring=[b]",
}
\lstdefinestyle{python}{
  language=Python,
  basicstyle=\ttfamily\footnotesize,
  backgroundcolor=\color{codebg},
  keywordstyle=\color{codekw}\bfseries,
  stringstyle=\color{codestr},
  commentstyle=\color{codecomment}\itshape,
  breaklines=true,
  frame=single,
  rulecolor=\color{black!20},
  numbers=left,
  numberstyle=\tiny\color{pynegray},
  numbersep=6pt,
  xleftmargin=1.5em,
  framexleftmargin=1.2em,
  showstringspaces=false,
  tabsize=4,
}
\lstset{style=python}

% Theorem environments
\theoremstyle{plain}
\newtheorem{theorem}{Theorem}[section]
\newtheorem{lemma}[theorem]{Lemma}
\newtheorem{proposition}[theorem]{Proposition}
\newtheorem{corollary}[theorem]{Corollary}
\theoremstyle{definition}
\newtheorem{definition}[theorem]{Definition}
\newtheorem{example}[theorem]{Example}
\newtheorem{remark}[theorem]{Remark}

% Shortcuts
\newcommand{\pyne}{\textsc{Pyne}}
\newcommand{\pine}{Pine Script\texttrademark}
\newcommand{\tv}{TradingView\textregistered}
\newcommand{\asdl}{\textsc{Asdl}}
\newcommand{\antlr}{\textsc{Antlr}4}
\newcommand{\numba}{\textsc{Numba}}
\newcommand{\lsp}{\textsc{Lsp}}
\newcommand{\ohlcv}{\textsc{Ohlcv}}
\newcommand{\udt}{\textsc{Udt}}
\newcommand{\jit}{\textsc{Jit}}
\newcommand{\api}{\textsc{Api}}
\newcommand{\cli}{\textsc{Cli}}
% AST acronym (text mode). Keep amsmath's math asterisk \ast intact.
\DeclareRobustCommand{\Ast}{\textsc{Ast}}
% Papers in this series write \ast{} for the AST acronym in prose.
% Redefine robustly; use \mathast for a binary asterisk in formulas.
\let\mathast\ast
\DeclareRobustCommand{\ast}{\textsc{Ast}}
\newcommand{\ir}{\textsc{Ir}}
\newcommand{\ta}{\texttt{ta}}
\newcommand{\code}[1]{\texttt{#1}}
\newcommand{\file}[1]{\texttt{#1}}
\newcommand{\eg}{e.g.}
\newcommand{\ie}{i.e.}
\newcommand{\etal}{\textit{et~al.}}
\newcommand{\resp}{respectively}

% Math helpers
\newcommand{\R}{\mathbb{R}}
\newcommand{\N}{\mathbb{N}}
\newcommand{\Z}{\mathbb{Z}}
% Text-mode NA token for prose; use $\namath$ inside formulas if needed.
\DeclareRobustCommand{\na}{\textsf{na}}
\newcommand{\namath}{\mathsf{na}}
\newcommand{\series}[1]{\mathbf{#1}}
\newcommand{\baridx}{i}
\newcommand{\bars}{n}
\newcommand{\period}{p}
\newcommand{\abserr}{\varepsilon_{\mathrm{abs}}}
\newcommand{\relerr}{\varepsilon_{\mathrm{rel}}}
\newcommand{\maxerr}{\varepsilon_{\mathrm{max}}}
\newcommand{\meanerr}{\varepsilon_{\mathrm{mean}}}

% Pipeline notation — use inside math: $\pipeline$
\newcommand{\pipeline}{\mathrm{src}\!\to\!\mathrm{lex}\!\to\!\mathrm{parse}\!\to\!\mathrm{AST}\!\to\!\mathrm{eval}}

% Disclaimer footnote helper
\newcommand{\trademarknote}{%
  \pine{} and \tv{} are trademarks of TradingView, Inc.
  \pyne{} is an independent, unofficial implementation and is not affiliated with,
  authorized by, sponsored by, or endorsed by TradingView, Inc.
}

% Figure path helper (papers live one level under .papers/)
\graphicspath{{figures/}{../common/figures/}}


\hypersetup{
  pdftitle={Efficient Series History and Incremental Technical Analysis for Bar-Loop Interpreters},
  pdfauthor={HOOX Research / PYNE Contributors},
  pdfkeywords={incremental computation, technical analysis, series buffers, ring buffers,
               memory-bounded interpreters, financial time series}
}

\title{%
  Efficient Series History and Incremental Technical Analysis\\
  for Bar-Loop Interpreters}
\author{%
  HOOX Research / \pyne{} Contributors\\
  \texttt{jango\_blockchained}\\
  Independent Open Source Research\\
  \texttt{contact@hoox.sh}}
\date{August 2026}

\begin{document}
\maketitle

\begin{abstract}
Bar-loop interpreters for series-oriented domain-specific languages evaluate
every statement once per time bar and materialize history via the operator
$x[n]$ (``$n$ bars back'').  Naive implementations store unbounded per-series
lists and recompute rolling technical-analysis (TA) windows from scratch on
each bar, inducing $O(n\cdot s)$ memory growth and $O(n\cdot p)$ work for
period-$p$ indicators over $n$ bars and $s$ series.
We present the series-history and incremental-TA stack shipped in \pyne{}, an
independent open toolchain for the \pine{} language~\citep{pynescript2026,tradingview-pine}.
The design combines (i)~a default-on series cap that trims chronological host
lists to $\max(\mathit{SERIES\_MAX},\mathit{max\_bars\_back})+\mathit{slack}$,
yielding $\sim 78\times$ long-run list-memory reduction at $20\,000$ bars;
(ii)~an opt-in chronological ring buffer (\texttt{ChronologicalSeriesBuffer})
with true $O(1)$ lookback under modular indexing;
(iii)~call-site incremental kernels for the hot \ta{} surface---including
SMA/EMA/RMA, RSI, Bollinger Bands, highest/lowest, KAMA, and StochRSI---with
kernel micro-benchmarks up to $\sim 717\times$ versus full-window rebuilds; and
(iv)~host-side parse/AST caching, lazy calendar materialization, and light plot
modes that cut multi-plot residual wall time.
On the interpret path at $2\,000$ bars we report median latencies of $15.2$\,ms
(minimal), $23.6$\,ms (\texttt{ta\_sma}), $153$\,ms (\texttt{ta\_combo}), and
$69$\,ms (\texttt{strategy\_ish}); cProfile attributes $\sim 75\%$ of residual
\texttt{ta\_combo} wall time to plot packing rather than TA kernels.
We formalize correctness conditions under bounded history, give recurrences and
stability notes for each kernel class, and situate the work in incremental
computation and streaming algorithms.
\end{abstract}

\paragraph{Keywords.}
incremental computation; technical analysis; series buffers; ring buffers;
memory-bounded interpreters; financial time series.

\paragraph{Disclaimer.}
\trademarknote{}

\tableofcontents
\newpage

%==============================================================================
\section{Introduction}
\label{sec:intro}
%==============================================================================

Financial charting languages such as \pine{}~\citep{tradingview-pine} expose a
\emph{series} data model in which every scalar expression is implicitly a time
series indexed by bar.  Evaluation proceeds as a \emph{bar loop}: the host
advances an \ohlcv{} cursor and re-executes the script body for each bar $i\in
\{0,\ldots,n-1\}$.  History access $x[n]$ denotes the value of $x$ that was
current $n$ bars earlier; qualifiers \code{var} and \code{varip} persist state
across bars and intra-bar ticks respectively; and the sentinel $\na$
propagates through arithmetic~\citep{tradingview-pine,pynescript2026}.

This model is attractive for domain users---indicators are written as if they
were spreadsheet formulas---but is expensive for a naive interpreter.
Historically, the dominant costs on \pyne{}'s interpret path were:
\begin{enumerate}[leftmargin=*]
  \item \textbf{Series history growth.}  Every named series and host \ohlcv{}
        field appended one sample per bar, so memory scaled as
        $\Theta(n\cdot s)$ for $s$ live series.
  \item \textbf{Full-window TA recompute.}  Indicators such as
        $\mathrm{SMA}(p)$, Bollinger Bands, and Kaufman adaptive moving
        average~\citep{kama1998kaufman} rebuilt period-$p$ windows from the
        entire prefix on each bar, producing $\Theta(n\cdot p)$ or even
        $\Theta(n^2)$ work for nested rebuilds.
  \item \textbf{Plot packing and call dispatch.}  Multi-plot scripts spent a
        majority of residual wall time in result capture and
        \texttt{visit\_Call} envelopes rather than pure arithmetic.
\end{enumerate}

\pyne{} addresses these costs with a layered runtime stack measured across
performance rounds~4--7~\citep{pynescript2026}.  Contributions of this paper:

\begin{itemize}[leftmargin=*]
  \item \textbf{Bounded series history} via default-on \code{PYNE\_SERIES\_CAP},
        with a correctness argument relating lookback depth, recursive
        incremental state, and out-of-range $\na$ semantics
        (Section~\ref{sec:cap}).
  \item \textbf{Chronological ring buffers}
        (\texttt{ChronologicalSeriesBuffer} /
        \code{PYNE\_SERIES\_RING}) enabling $O(1)$ Pine-offset lookback and a
        dual-write migration path off newest-first deques
        (Section~\ref{sec:ring}).
  \item \textbf{A generic incremental TA framework} of call-site state tuples
        $\sigma$ and updates $\sigma'=\mathrm{update}(\sigma,x_i)$, covering
        sliding sums, exponential smoothers, Wilder-class oscillators, variance
        bands, amortized deques, and adaptive kernels
        (Section~\ref{sec:inc}).
  \item \textbf{Interpreter and multi-run host optimizations}: type-keyed
        dispatch, light plots, lazy calendar, and a SHA-256 LRU \ast{} cache
        with call-site scrubbing on hit (Sections~\ref{sec:hot}--\ref{sec:parse}).
  \item \textbf{Empirical evaluation} of memory, kernel micros, and end-to-end
        latency, including a residual bottleneck map that isolates plot packing
        as $\sim 75\%$ of multi-plot interpret wall time
        (Section~\ref{sec:eval}).
\end{itemize}

\paragraph{Scope.}
We focus on the \emph{interpret} (AST-walking) path and on the series/TA
subsystem shared with the Numba compile path~\citep{lam2015numba}, which always
emits incremental kernels for hot TA.  Grammar, strategy broker semantics, and
LSP tooling are treated elsewhere in the \pyne{} paper series.

%==============================================================================
\section{Background}
\label{sec:bg}
%==============================================================================

\subsection{Pine series semantics}
\label{sec:bg-series}

\begin{definition}[Series and history operator]
A series $\series{x}$ is a sequence $(x_0,x_1,\ldots,x_{i})$ of values observed
up to the current bar index $i$.  The history operator is
\begin{equation}
  x[n] \;=\;
  \begin{cases}
    x_{i-n} & \text{if } 0\le n\le i \text{ and } x_{i-n}\ne\na,\\
    \na     & \text{otherwise (including } n<0\text{).}
  \end{cases}
\end{equation}
Pine indices are \emph{most-recent-first} in the programmer model ($x[0]$ is
current), while host chronological lists store oldest-first so that physical
index $-(n+1)$ realizes lookback $n$~\citep{tradingview-pine}.
\end{definition}

\begin{definition}[Persistence qualifiers]
The qualifier \code{var} initializes a declaration site once and carries the
value across subsequent bars.  \code{varip} additionally carries across ticks
within an incomplete bar.  Bare assignment re-evaluates every bar.
\end{definition}

\begin{definition}[$\na$ propagation]
Arithmetic and many builtins propagate $\na$: if any operand is $\na$, the
result is $\na$ (with Pine-specific exceptions for functions that treat
missing samples specially).  Out-of-range history must yield $\na$, never a
silent zero~\citep{pynescript2026}.
\end{definition}

These rules interact with memory bounding: truncating history deeper than any
live lookback must not invent samples, and recursive smoothers must either
retain enough raw history \emph{or} retain closed-form incremental state
(Section~\ref{sec:cap}).

\subsection{Bar-loop evaluation model}
\label{sec:bg-barloop}

Figure~\ref{fig:architecture} sketches the interpret architecture.  For each
bar $i$:
\begin{enumerate}[leftmargin=*]
  \item The host advances \ohlcv{} series (\code{open}, \code{high},
        \code{low}, \code{close}, \code{volume}, \code{time}) and optional
        script-local series.
  \item The evaluator walks the script \ast{} via a visitor, resolving names
        against a context map and call-site TA state.
  \item Builtins (notably \ta{}) either update incremental kernels or fall
        back to full-window helpers when \code{PYNE\_TA\_INCREMENTAL}${}=0$.
  \item Plot/strategy side effects append into result registries for packing
        into columnar outputs.
\end{enumerate}

The compile path lowers the same semantics to NumPy/Numba kernels that process
entire arrays in one call~\citep{harris2020numpy,lam2015numba}; that path always
uses incremental kernels for hot TA and is not gated by the interpret flag.

\begin{figure}[t]
\centering
\begin{tikzpicture}[
  font=\small,
  box/.style={draw, rounded corners=2pt, align=center, inner sep=6pt,
              minimum height=1.1cm, minimum width=2.4cm},
  arr/.style={-{Latex}, thick},
  host/.style={box, fill=pyneblue!12},
  series/.style={box, fill=pyneorange!15},
  eval/.style={box, fill=pynegreen!12},
  plot/.style={box, fill=pynegray!15},
]
  \node[host] (ohlcv) {OHLCV\\host cursor};
  \node[series, right=1.1cm of ohlcv] (store) {Series store\\lists / ring};
  \node[eval, right=1.1cm of store] (ast) {AST visitor\\per bar $i$};
  \node[eval, below=1.0cm of ast] (ta) {TA kernels\\state $\sigma$};
  \node[plot, right=1.1cm of ast] (plots) {Plot / strategy\\registries};
  \node[plot, below=1.0cm of plots] (pack) {Columnar pack\\result dict};
  \node[host, below=1.0cm of ohlcv] (bar) {Bar loop\\$i=0..n-1$};

  \draw[arr] (bar) -- (ohlcv);
  \draw[arr] (ohlcv) -- (store);
  \draw[arr] (store) -- (ast);
  \draw[arr] (ast) -- (plots);
  \draw[arr] (ast) -- (ta);
  \draw[arr] (ta) -- (ast);
  \draw[arr] (plots) -- (pack);
  \draw[arr, dashed] (store) -- (ta);

  \node[draw, dashed, rounded corners=4pt, fit=(ohlcv)(store)(bar),
        inner sep=8pt, label=above:{\footnotesize Host}] {};
  \node[draw, dashed, rounded corners=4pt, fit=(ast)(ta),
        inner sep=8pt, label=above:{\footnotesize Interpreter}] {};
  \node[draw, dashed, rounded corners=4pt, fit=(plots)(pack),
        inner sep=8pt, label=above:{\footnotesize Effects}] {};
\end{tikzpicture}
\caption{Bar-loop interpret architecture with series store, call-site TA
state, and plot packing.  The host advances OHLCV once per bar; the visitor
re-executes the script body; incremental kernels read/write state $\sigma$
keyed by call site.}
\label{fig:architecture}
\end{figure}

\subsection{Cost model of interpretation}
\label{sec:bg-cost}

Let $n$ be the number of bars, $s$ the number of live series (host + local),
$c$ the number of call nodes executed per bar, $p$ a typical TA period, and
$q$ the number of \code{plot} (or similar) emissions per bar.  A crude cost
model for a naive interpreter is
\begin{equation}
\label{eq:cost-naive}
  T_{\mathrm{naive}}(n)
  \;=\;
  \underbrace{O(n\cdot c\cdot d)}_{\text{AST dispatch depth }d}
  \;+\;
  \underbrace{O\!\big(n\cdot \textstyle\sum_k p_k\big)}_{\text{full-window TA}}
  \;+\;
  \underbrace{O(n\cdot q\cdot w)}_{\text{plot capture width }w}
  \;+\;
  \underbrace{O(n\cdot s)}_{\text{series append}}.
\end{equation}
Memory is $M_{\mathrm{naive}}=\Theta(n\cdot s)$ plus plot buffers
$\Theta(n\cdot q)$.

With series caps of capacity $K$, incremental TA of $O(1)$ or $O(p)$ amortized
per call, and light plots that skip columnar capture, the improved model is
\begin{equation}
\label{eq:cost-opt}
  T_{\mathrm{opt}}(n)
  \;=\;
  O(n\cdot c\cdot d')
  \;+\;
  O\!\big(n\cdot \textstyle\sum_k u_k\big)
  \;+\;
  O(n\cdot q\cdot w')
  \;+\;
  O(n\cdot s),
\end{equation}
where $u_k\in\{1,p_k\}$ is the per-bar update cost of kernel $k$, $d'\le d$
after dispatch tuning, $w'\ll w$ under light plots, and
$M_{\mathrm{opt}}=\Theta(K\cdot s)+\Theta(|\sigma|)$ with $|\sigma|$ the
aggregate incremental state (typically $O(p)$ per call site).

\paragraph{Dominant historical costs.}
Profiling rounds identified four hotspots: series list growth, TA full
recompute, plot packing, and the \texttt{visit\_Call} envelope.  Kernel-only
wins therefore do not automatically dominate end-to-end multi-plot scripts
(Section~\ref{sec:eval}).

%==============================================================================
\section{Problem Formalization}
\label{sec:problem}
%==============================================================================

\subsection{Memory: unbounded series growth}

\begin{proposition}[Linear series memory]
\label{prop:mem-linear}
Under uncapped chronological lists, if a script maintains $s$ series each of
length $n$ after $n$ bars, then pointer storage alone is
$\Theta(n\cdot s)$ words.  For $n=20\,000$, $s=6$ OHLCV-style lists, this is
$\sim 9.6\cdot 10^5$ pointers ($\sim 7.7$\,MiB on 64-bit), exclusive of object
payloads.
\end{proposition}

Many Pine scripts also allocate local series for intermediate expressions;
$s$ can reach dozens to hundreds on complex strategies, so memory pressure is
not limited to OHLCV.

\subsection{Time: full-window TA recompute}

\begin{definition}[Full-window recompute]
A kernel $f_p$ is a \emph{full-window recompute} if at bar $i$ it scans
$\{x_{i-p+1},\ldots,x_i\}$ (or the entire prefix) without reusing work from
bar $i-1$, costing $\Omega(p)$ or $\Omega(i)$ arithmetic per bar.
\end{definition}

\begin{proposition}[Quadratic nested rebuild]
\label{prop:quad}
If each bar $i$ rebuilds an indicator by scanning a prefix of length $i$
(e.g., nested KAMA or StochRSI implementations that re-seed from bar~$0$),
total work over $n$ bars is $\Theta(n^2)$.
\end{proposition}

Even ``only'' $\Theta(n\cdot p)$ is painful for $n=5\,000$ and $p=200$ when
dozens of indicators share a script.  Time-series textbooks emphasize
online/streaming updates for exactly this reason~\citep{box1970time,tsay2010analysis}.

\subsection{Correctness constraints}

Any optimization must preserve:
\begin{enumerate}[leftmargin=*]
  \item \textbf{Lookback identity:} for all $n\le K$ (cap), $x[n]$ equals the
        uncapped value when the sample exists.
  \item \textbf{$\na$ discipline:} OOB and missing samples are $\na$, never $0$.
  \item \textbf{TA parity:} incremental kernels match full-window oracles within
        floating-point noise~\citep{higham2002accuracy,ieee754} under default
        flags.
  \item \textbf{Call-site isolation:} two syntactically distinct calls to
        \code{ta.sma(close,14)} maintain independent state.
\end{enumerate}

%==============================================================================
\section{Series Cap Design}
\label{sec:cap}
%==============================================================================

\subsection{Policy}

\pyne{} formalizes host list capping behind \code{PYNE\_SERIES\_CAP}
(default \textbf{on}).  Let
\begin{equation}
  K \;=\; \max\!\big(K_0,\; \mathit{max\_bars\_back}(\mathrm{src})\big),
\end{equation}
where $K_0=\mathit{SERIES\_MAX}$ defaults to $256$ (overridable via
\code{PYNE\_SERIES\_MAX}) and $\mathit{max\_bars\_back}$ is parsed from the
script source when present.  Chronological lists are grown to $K+\Delta$ with
slack $\Delta=64$, then trimmed in place with \texttt{del lst[:drop]} back to
$K$.  Slack amortizes the $O(K)$ cost of prefix deletion so trims are not
performed every bar.

PineSeries wrappers retain a history floor of $1\,000$ samples (raised when
$K$ is larger) so \code{close[n]} for $n\le 999$ remains available even when
list caps are tight---matching pre-cap host defaults.

\subsection{Correctness argument for bounded history}

\begin{proposition}[Window kernels under cap]
\label{prop:window-cap}
Let $f_p$ depend only on the last $p$ finite samples of $\series{x}$.  If
$p\le K$ and the cap retains the newest $K$ samples, then $f_p$ at every bar
with $i\ge p-1$ equals the uncapped result.
\end{proposition}

\begin{proof}[Proof sketch]
At bar $i$, the uncapped window is $\{x_{i-p+1},\ldots,x_i\}$.  The cap
retains $\{x_{i-K+1},\ldots,x_i\}$ (or the full prefix if shorter).  Since
$p\le K$, the window is a suffix of the retained segment.
\end{proof}

\begin{proposition}[Recursive kernels with incremental state]
\label{prop:rec-cap}
Let $y_i=g(y_{i-1},x_i)$ be a first-order recurrence (EMA, RMA, \ldots) with
call-site state storing $y_{i-1}$.  After warm-up, $y_i$ is independent of raw
samples older than $1$, hence independent of $K$ provided state is not reset.
\end{proposition}

\begin{remark}[Full-recompute residual]
If \code{PYNE\_TA\_INCREMENTAL}${}=0$, recursive kernels re-seed from the
retained list prefix.  When $n\gg K$, the re-seeded trajectory can diverge from
an uncapped full recompute.  Default-on incremental mode closes this residual
for production hosts; the flag remains for debugging and oracle tests.
\end{remark}

\begin{proposition}[OOB is $\na$]
\label{prop:oob-na}
For lookback $n\ge$ retained length, $x[n]=\na$.  Implementations must not
coerce missing history to $0$, which would bias sums, ranges, and boolean
guards.
\end{proposition}

\subsection{When lookback exceeds the cap}

If a script requests \code{close[500]} while $K=256$, the host returns $\na$
unless $\mathit{max\_bars\_back}$ (or \code{PYNE\_SERIES\_MAX}) raises $K$.
This matches Pine's resource model: deep history is an explicit request, not
an implicit infinite tape.  Scripts that need deep history should set
\code{max\_bars\_back} accordingly; the cap then becomes a \emph{floor} rather
than a truncation surprise.

\subsection{Structural memory win}

Simulating $20\,000$ bars $\times$ $6$ OHLCV-style lists with $K=256$,
$\Delta=64$:

\begin{center}
\begin{tabular}{lrrr}
\toprule
Mode & Final list len & $\sim$pointer bytes (6 lists) & Append wall \\
\midrule
Uncapped & $20\,000$ & $\sim 960\,000$ & $0.020$\,s \\
Capped   & $305$     & $\sim 14\,640$  & $0.021$\,s \\
\bottomrule
\end{tabular}
\end{center}

The memory ratio is $\sim 78\times$ at $20$k bars and scales as
$n/K$ asymptotically.  Trim CPU is negligible relative to AST and TA work.

\begin{figure}[t]
\centering
\begin{tikzpicture}
\begin{axis}[
  width=0.92\linewidth,
  height=6.2cm,
  xlabel={Bar count $n$},
  ylabel={Relative series memory (a.u.)},
  legend style={at={(0.02,0.98)}, anchor=north west, font=\small},
  grid=major,
  ymin=0,
  xmax=20000,
  cycle list name=color list,
]
  \addplot[thick, color=pynegray, mark=none] coordinates {
    (0,0) (2500,2500) (5000,5000) (7500,7500) (10000,10000)
    (12500,12500) (15000,15000) (17500,17500) (20000,20000)
  };
  \addlegendentry{Uncapped $\Theta(n)$}
  \addplot[thick, color=pyneorange, mark=none] coordinates {
    (0,0) (256,256) (320,320) (1000,320) (5000,320)
    (10000,320) (15000,320) (20000,320)
  };
  \addlegendentry{Capped $K{+}{\Delta}$ ($K{=}256$)}
  \addplot[thick, color=pyneblue, dashed, mark=none] coordinates {
    (0,0) (256,256) (1000,256) (5000,256) (10000,256)
    (15000,256) (20000,256)
  };
  \addlegendentry{Ring $\mathrm{maxlen}{=}K$}
\end{axis}
\end{tikzpicture}
\caption{Series growth versus capped lists and fixed-capacity ring buffers.
Uncapped memory grows linearly with bar count; the default series cap plateaus
near $K+\Delta$; a modular ring plateaus exactly at $K$.}
\label{fig:memory}
\end{figure}

%==============================================================================
\section{Ring Buffer Series}
\label{sec:ring}
%==============================================================================

\subsection{\texttt{ChronologicalSeriesBuffer}}

The Phase~2.2 design introduces a chronological (oldest-first) buffer with
modular ring capacity~\citep{pynescript2026}:

\begin{align}
  \mathrm{append}(v) &\quad O(1), \\
  \mathrm{lookback}(n) &\quad O(1)\text{ via physical index }\mathit{end}-n, \\
  \mathrm{series}[n] &\quad \na\text{ on negative / OOB / missing}.
\end{align}

When $\mathit{maxlen}=N$ is set, append overwrites the oldest slot after the
ring fills---true $O(1)$ memory and time without prefix \texttt{del}.  When
$\mathit{maxlen}$ is \code{None}, the buffer grows as a plain list (still $O(1)$
end-index lookback).

\paragraph{Why not \texttt{deque} alone?}
CPython \texttt{collections.deque} offers $O(1)$ append/pop at both ends, but
\emph{middle} indexing is $O(n)$.  Modular list indexing is $O(1)$ for
arbitrary lookback depth within the window---important for scripts that probe
\code{close[k]} for moderate $k$ on every bar.

\subsection{Dual-write transition}

Production hosts historically maintained \emph{two} series representations:
\begin{itemize}[leftmargin=*]
  \item \texttt{PineSeries}: newest-first deque (\texttt{appendleft}), used by
        context OHLCV and \code{series[n]}.
  \item \texttt{current\_series} lists: chronological, used by \ta{} via
        \texttt{\_as\_series}.
\end{itemize}
Materializing a newest-first deque into chronological order required
\texttt{list(reversed(...))}, a residual tax on full-recompute paths.

The migration plan is staged:
\begin{enumerate}[leftmargin=*]
  \item \textbf{Phase A (shipped):} factory \texttt{make\_pine\_series()}
        returns \texttt{RingPineSeries} when \code{PYNE\_SERIES\_RING}${}=1$,
        else classic \texttt{PineSeries}.  Default is \textbf{off} until the
        corpus is green under ring mode.
  \item \textbf{Phase B:} zero-copy \texttt{\_as\_series} for objects marked
        \code{chrono\_order=True}.
  \item \textbf{Phase C:} single buffer shared by context and
        \texttt{current\_series}, dropping dual update.
\end{enumerate}

\texttt{RingPineSeries} exposes a \texttt{NewestFirstHistoryView} so existing
duck-typed code (\code{history[0]} == current) continues to work during the
transition.  Ring \code{maxlen} tracks
\texttt{pineseries\_history\_length(...)} so it composes with the series cap
rather than fighting it.

%==============================================================================
\section{Incremental TA Framework}
\label{sec:inc}
%==============================================================================

\subsection{Generic pattern}
\label{sec:inc-pattern}

\begin{definition}[Incremental kernel]
An incremental kernel is a pair $(\sigma_0,U)$ where $\sigma_0$ is initial
state and $U(\sigma,x_i)=(\sigma',y_i)$ produces updated state and output.
State is stored in a call-site map keyed by (builtin name, syntactic slot,
parameters).
\end{definition}

\begin{algorithm}[t]
\caption{Generic bar-mode incremental update}
\label{alg:generic}
\begin{algorithmic}[1]
\Require call-site map $S$, key $k$, input $x_i$, update $U$, init $\sigma_0$
\If{$k\notin S$}
  \State $S[k]\gets\sigma_0$
\EndIf
\State $(\sigma', y_i)\gets U(S[k], x_i)$
\State $S[k]\gets\sigma'$
\State \Return $y_i$ \Comment{may be $\na$ during warm-up}
\end{algorithmic}
\end{algorithm}

Flag \code{PYNE\_TA\_INCREMENTAL} (default \textbf{on}) selects incremental
updates; setting it to $0$ forces full-window oracles for differential testing.
The compile path always emits incremental Numba kernels for hot TA regardless
of this flag~\citep{lam2015numba,pynescript2026}.

Figure~\ref{fig:inc-schematic} contrasts full-window scans with incremental
updates.

\begin{figure}[t]
\centering
\begin{tikzpicture}[
  font=\small,
  bar/.style={draw, minimum width=0.42cm, minimum height=0.7cm, inner sep=0},
  arr/.style={-{Latex}, thick},
]
  % Full window
  \node[anchor=west] at (0,2.4) {\textbf{Full-window recompute at bar $i$}};
  \foreach \x in {0,...,11} {
    \node[bar, fill=pynegray!20] (f\x) at (0.5*\x, 1.6) {};
  }
  \foreach \x in {7,...,11} {
    \node[bar, fill=pyneorange!50] at (0.5*\x, 1.6) {};
  }
  \draw[decorate, decoration={brace, amplitude=5pt, mirror}, thick]
    (3.35,1.15) -- (5.65,1.15)
    node[midway, below=6pt] {scan $p$ samples, $O(p)$};
  \node at (6.6,1.6) {$y_i$};

  % Incremental
  \node[anchor=west] at (0,0.3) {\textbf{Incremental update}};
  \node[draw, rounded corners, fill=pyneblue!12, minimum width=2.2cm,
        minimum height=1.0cm] (st) at (1.4,-0.6) {state $\sigma_{i-1}$};
  \node[draw, rounded corners, fill=pynegreen!15, minimum width=1.4cm,
        minimum height=1.0cm] (xi) at (3.8,-0.6) {$x_i$};
  \node[draw, rounded corners, fill=pyneorange!20, minimum width=2.2cm,
        minimum height=1.0cm] (st2) at (6.4,-0.6) {$\sigma_i,\; y_i$};
  \draw[arr] (st) -- (st2);
  \draw[arr] (xi) -- (st2);
  \node[below=0.15cm of st2, font=\footnotesize] {$O(1)$ or amortized $O(1)$};
\end{tikzpicture}
\caption{Incremental versus full-window TA.  Full recompute re-scans the
period window on every bar; incremental kernels carry a state tuple and apply
a constant-time (or amortized constant-time) update.}
\label{fig:inc-schematic}
\end{figure}

\subsection{Case study: SMA (sliding sum)}
\label{sec:inc-sma}

Let $p\ge 1$.  Maintain a deque window $W$ of the last $p$ samples and a
running sum $\Sigma$:
\begin{align}
  \Sigma_i &= \Sigma_{i-1} + x_i - x_{i-p}
    \quad\text{(when $|W|=p$)}, \\
  y_i &= \Sigma_i / p.
\end{align}
Warm-up returns $\na$ until $p$ finite samples accumulate.  Samples that are
$\na$ are typically excluded from the count or propagate $\na$ per Pine rules.

\begin{algorithm}[t]
\caption{Incremental SMA with running sum}
\label{alg:sma}
\begin{algorithmic}[1]
\Require state $(\mathit{window},\Sigma,\mathit{count})$, input $x$, period $p$
\If{$x$ is $\na$}
  \State \Return $(\text{state}, \na)$ \Comment{policy: skip or propagate}
\EndIf
\State append $x$ to $\mathit{window}$; $\Sigma\gets\Sigma+x$; $\mathit{count}{+}{+}$
\If{$\mathit{count}>p$}
  \State $u\gets$ popleft($\mathit{window}$); $\Sigma\gets\Sigma-u$; $\mathit{count}{-}{+}$
\EndIf
\If{$\mathit{count}<p$}
  \State \Return $(\text{state}, \na)$
\EndIf
\State \Return $(\text{state}, \Sigma/p)$
\end{algorithmic}
\end{algorithm}

Complexity: $O(1)$ arithmetic and deque operations per bar after warm-up.

\subsection{Case study: EMA and RMA (Wilder)}
\label{sec:inc-ema}

The exponential moving average with smoothing $\alpha\in(0,1]$ is
\begin{equation}
\label{eq:ema}
  y_i = \alpha x_i + (1-\alpha)\, y_{i-1},
\end{equation}
seeded by SMA of the first $p$ samples for Pine-compatible EMA with
$\alpha=2/(p+1)$.  Wilder's RMA~\citep{wilder1978new} uses
$\alpha=1/p$ (equivalently $y_i = y_{i-1} + (x_i-y_{i-1})/p$).

State is a single scalar $y_{i-1}$ plus a warm-up counter---$O(1)$ memory and
time.  Numerical stability:~\eqref{eq:ema} is a contractive affine map; error
is not amplified~\citep{higham2002accuracy}.  Catastrophic cancellation is not
an issue for the recurrence itself, though very large $|x_i|$ can still lose
precision in IEEE-754 binary64~\citep{ieee754}.

\subsection{Case study: RSI}
\label{sec:inc-rsi}

Relative Strength Index~\citep{wilder1978new} maintains average gain $G$ and
average loss $L$ via RMA:
\begin{align}
  u_i &= \max(x_i-x_{i-1}, 0), &
  d_i &= \max(x_{i-1}-x_i, 0), \\
  G_i &= \mathrm{RMA}_p(u)_i, &
  L_i &= \mathrm{RMA}_p(d)_i, \\
  \mathrm{RS}_i &= G_i/L_i, &
  \mathrm{RSI}_i &= 100 - 100/(1+\mathrm{RS}_i).
\end{align}
When $L_i=0$, RSI is defined as $100$ if $G_i>0$, else typically $0$ or $\na$
per implementation policy.  Incremental form stores
$(G_{i-1},L_{i-1},x_{i-1})$.

\subsection{Case study: Bollinger Bands (mean + variance)}
\label{sec:inc-bb}

Bollinger Bands~\citep{bollinger2001} are
\begin{equation}
  \mathrm{mid}_i=\mathrm{SMA}_p(x)_i,\quad
  \sigma_i=\mathrm{stdev}_p(x)_i,\quad
  \mathrm{upper/lower}_i=\mathrm{mid}_i \pm k\sigma_i.
\end{equation}
\pyne{} maintains sliding sum $\Sigma$ and sum of squares $Q$:
\begin{align}
  \Sigma_i &= \Sigma_{i-1}+x_i-x_{i-p}, \\
  Q_i &= Q_{i-1}+x_i^2-x_{i-p}^2, \\
  \sigma_i^2 &= Q_i/p - (\Sigma_i/p)^2
  \quad\text{(population form; sample form uses $p-1$)}.
\end{align}
Alternatively, Welford's online algorithm updates mean and $M_2$ with better
numerical properties for one-pass streams~\citep{higham2002accuracy,press2007numerical};
for fixed-length sliding windows, a compensated sliding variance or periodic
full refresh reduces cancellation when $Q$ and $(\Sigma/p)^2$ are large and
close.

Round~7 ships a dedicated \texttt{\_bb\_inc\_update} that nests SMA+stdev
slots and last-sample hygiene so Volatility paths avoid reverse materialization.
Kernel micro-benchmarks report $\sim 217\times$ versus naive full
\texttt{\_sma}+\texttt{\_stdev} rebuilds at $n=2\,000$
(Section~\ref{sec:eval}).

\subsection{Case study: highest / lowest (monotonic deques)}
\label{sec:inc-hl}

Range extrema over a sliding window admit classic amortized $O(1)$ updates via
monotonic deques~\citep{aho2006dragon}:

\begin{algorithm}[t]
\caption{Sliding window maximum (highest)}
\label{alg:highest}
\begin{algorithmic}[1]
\Require deque $D$ of indices (decreasing values), input $x_i$, period $p$
\While{$D$ nonempty and $x_{D.\mathrm{back}}\le x_i$}
  \State pop back $D$
\EndWhile
\State push $i$ to back of $D$
\While{$D.\mathrm{front}\le i-p$}
  \State pop front $D$
\EndWhile
\State \Return $x_{D.\mathrm{front}}$ \Comment{$\na$ if window incomplete}
\end{algorithmic}
\end{algorithm}

Each index is pushed and popped at most once, so amortized cost is $O(1)$ per
bar.  Lowest is symmetric (increasing deque).

\subsection{Case study: KAMA}
\label{sec:inc-kama}

Kaufman's adaptive moving average~\citep{kama1998kaufman} computes an
efficiency ratio from change over noise:
\begin{align}
  \mathrm{ER}_i &= \frac{|x_i-x_{i-L}|}{\sum_{j=0}^{L-1}|x_{i-j}-x_{i-j-1}|}, \\
  \mathrm{SC}_i &= \big(\mathrm{ER}_i\cdot(\mathrm{fast}-\mathrm{slow})+\mathrm{slow}\big)^2, \\
  y_i &= y_{i-1} + \mathrm{SC}_i\cdot(x_i-y_{i-1}).
\end{align}
A naive implementation rebuilds the noise sum from scratch each bar
($O(L)$ or worse if re-seeded from bar~$0$).  The incremental kernel keeps a
price ring of length $L{+}1$ and a running absolute-delta sum, achieving
$O(1)$ per bar after seeding.  Micro-benchmarks: $\sim 121\times$ at
$n=2\,000$, length $10$.

\subsection{Case study: StochRSI}
\label{sec:inc-stochrsi}

Stochastic RSI applies a stochastic oscillator to an RSI series.  Nested
full rebuilds yield near-quadratic cost.  The incremental form maintains:
\begin{itemize}[leftmargin=*]
  \item an RSI window (or Wilder RSI state, depending on oracle),
  \item a stochastic ring over the last $\mathit{stoch\_len}$ RSI values,
  \item optional signal smoothing in call-site state
        ($0.33\,x+0.67\,\mathrm{prev}$ in \pyne{}'s simple-path oracle).
\end{itemize}
Round~7 reports $\sim 717\times$ kernel speedup for
\texttt{stochrsi(14,14)} at $n=2\,000$.  Note: the interpret oracle for this
path uses simple average gain/loss RSI rather than Wilder \code{ta.rsi}; the
incremental path matches that oracle bit-for-bit, not necessarily live \tv{}
StochRSI-on-RMA~\citep{pynescript2026}.

\subsection{Additional kernels}

Across rounds~1--7 the incremental surface grew to include (non-exhaustively):
WMA, VWMA, TR/ATR, change, stoch \%K, cum, VWAP, MOM, SWMA, barssince,
highestbars/lowestbars, linreg, CCI, TSI, ROC, WPR, dev, HMA, rising/falling,
median, percentrank, CMO, and compile-side ports of the same recurrences.
Typical gains range from $\sim 1.2\times$ (already $O(p)$ windows with better
hygiene) to two--three orders of magnitude for former full-prefix rebuilds.

\subsection{Numerical stability notes}
\label{sec:inc-stability}

\begin{itemize}[leftmargin=*]
  \item \textbf{Sliding variance:} prefer compensated updates or periodic
        resync when $|x|$ is large~\citep{higham2002accuracy}.
  \item \textbf{EMA/RMA:} contractive; warm-up seed choice dominates early bars.
  \item \textbf{Ratios (RSI, ER):} guard zero denominators explicitly; do not
        rely on IEEE infinities for Pine $\na$ semantics.
  \item \textbf{Parity budget:} compile vs.\ interpret max abs error
        $\sim 10^{-11}$ on standard kernels (float noise), $0$ on pure
        integer-path RSI/TSI variants in Numba goldens.
\end{itemize}

%==============================================================================
\section{Interpreter Hot Paths}
\label{sec:hot}
%==============================================================================

\subsection{Dispatch}

Expression evaluation is visitor-based~\citep{gamma1994design,nystrom2021crafting}.
Hot improvements include type-keyed dispatch tables (avoiding long
\texttt{isinstance} chains), operator maps for binary ops, and a scalar fast
path that skips series coercion when both operands are plain floats.  Mixed
expression micro-benchmarks improved $\sim 2$--$7\times$ across rounds.

The cumulative cProfile view on multi-plot TA scripts still shows
\texttt{visit\_Call} as a large envelope: thousands of calls per bar for
nested builtins, each paying Python function and argument-binding overhead.
This is structural to AST interpretation; the compile path eliminates it.

\subsection{Plot packing}

For \texttt{ta\_combo}-class scripts (multiple TA values, $\sim 8$ plots),
variant isolation shows:
\begin{center}
\begin{tabular}{lr}
\toprule
Variant (interpret, $n=2\,000$) & med wall \\
\midrule
\texttt{ta\_combo} with $8\times$\code{plot(...)} & $\sim 480$\,ms (profile host) \\
Same TA, \textbf{no plots} & $\sim 102$\,ms \\
\bottomrule
\end{tabular}
\end{center}
Thus plot path + host result packing account for $\sim 75$--$80\%$ of
multi-plot wall time once TA is incremental.  cProfile ranks
\texttt{\_builtin\_plot} / \texttt{\_capture\_plot} immediately after the
visitor envelope.

\subsection{Light plots and lazy calendar}

\begin{itemize}[leftmargin=*]
  \item \textbf{Lazy calendar} (always on): defer expensive calendar/timezone
        materialization until a plot or time builtin requires it.
  \item \textbf{\code{PYNE\_LIGHT\_PLOTS}:} skip columnar capture and registry
        bookkeeping when the consumer only needs OK/fail (corpus runs, smoke
        tests).  Multi-plot scripts see roughly $10$--$20\%$ wall reduction.
\end{itemize}

These are host-level levers: they do not change TA math, but they reclaim the
dominant residual identified by profiling.

\begin{figure}[t]
\centering
\begin{tikzpicture}
\begin{axis}[
  width=0.75\linewidth,
  height=7.0cm,
  ybar,
  bar width=18pt,
  ylabel={Kernel speedup ($\times$ vs full recompute)},
  symbolic x coords={kama, bb, stochrsi, vwap, kc, cmo},
  xtick=data,
  nodes near coords,
  nodes near coords align={vertical},
  every node near coord/.append style={font=\scriptsize},
  ymin=0, ymax=800,
  enlarge x limits=0.15,
  grid=major,
  tick label style={font=\small},
]
  \addplot[fill=pyneorange!70, draw=pyneorange!80!black] coordinates {
    (kama,121) (bb,217) (stochrsi,717) (vwap,217) (kc,78) (cmo,1.2)
  };
\end{axis}
\end{tikzpicture}
\caption{Kernel micro-benchmark speedups (incremental vs.\ full-window
rebuild) at $n\approx 2\,000$ bars.  Values from performance rounds~4--7:
KAMA $\sim 121\times$, Bollinger $\sim 217\times$, StochRSI $\sim 717\times$,
VWAP $\sim 217\times$, KC $\sim 78\times$, CMO $\sim 1.2\times$ (already
near-$O(p)$).}
\label{fig:speedups}
\end{figure}

\begin{figure}[t]
\centering
\begin{tikzpicture}
\begin{axis}[
  width=0.72\linewidth,
  height=7.0cm,
  x=3.2cm,
  y=3.2cm,
  xmin=-1.3, xmax=1.3,
  ymin=-1.3, ymax=1.3,
  axis equal,
  ticks=none,
  axis lines=none,
  clip=false,
]
  % pie slices via partial disks (approx)
  % plot packing ~75%, TA ~15%, other ~10%
  \fill[pyneorange!70] (0,0) -- (90:1) arc (90:-180:1) -- cycle;
  \fill[pyneblue!60] (0,0) -- (-180:1) arc (-180:-234:1) -- cycle;
  \fill[pynegray!50] (0,0) -- (-234:1) arc (-234:-270:1) -- cycle;
  \draw[thick] (0,0) circle (1);
  \node[font=\small] at (0.15,-0.35) {\textbf{Plot packing}};
  \node[font=\small] at (0.15,-0.55) {$\sim 75\%$};
  \node[font=\small] at (-0.55,0.55) {\textbf{TA}};
  \node[font=\small] at (-0.55,0.35) {$\sim 15\%$};
  \node[font=\small] at (0.55,0.75) {\textbf{Other}};
  \node[font=\small] at (0.55,0.55) {$\sim 10\%$};
\end{axis}
\end{tikzpicture}
\caption{Approximate latency breakdown for interpret \texttt{ta\_combo} after
incremental TA: plot packing dominates residual wall time; pure TA kernels are
a minority slice; ``Other'' covers dispatch, series updates, and host
bookkeeping.}
\label{fig:pie}
\end{figure}

%==============================================================================
\section{Parse/AST Caching Interaction with Runtime Multi-Run}
\label{sec:parse}
%==============================================================================

Product hosts often parse once and evaluate many times (parameter sweeps,
watchlists, recharts).  \pyne{} maintains a process-level LRU cache keyed by
SHA-256 of source text.  Multi-parse of identical source improves by
$\sim 1000\times$ on cache hit.

\paragraph{Call-site scrub on hit.}
AST nodes may accumulate mutable annotations (resolved call sites, cached
locations).  Returning the same \texttt{Script} object on a cache hit without
scrubbing caused empty-plot poisoning on the second \texttt{Runtime.run} of
identical source: call-site maps and plot registries interacted with stale
node state.  The fix clears Pine call-site annotations on cache hit
(\texttt{\_scrub\_pine\_call\_sites}) and pairs multi-run goldens with
\texttt{clear\_parse\_cache()} isolation where tests construct independent
Runtime instances~\citep{pynescript2026}.

\paragraph{Interaction with series/TA state.}
Parse caching is orthogonal to series caps and TA state: those live in the
evaluator/runtime instance.  However, multi-run correctness requires that
(1)~AST mutation not leak across runs, and (2)~evaluator call-site maps reset
or key correctly when scripts are re-bound.  Warm product paths prefer the
compile cache (near-zero warm compile, sub-millisecond runs for pure TA) when
scripts are compilation-eligible.

%==============================================================================
\section{Evaluation}
\label{sec:eval}
%==============================================================================

\subsection{Methodology}

Unless noted, numbers are from \pyne{} performance rounds~4--7 on Linux x86\_64
CPython, using \texttt{scripts/bench\_pipeline.py} and kernel micro-benches
(median of small iteration counts after warmup).  Absolute milliseconds vary
by host; we report the Round~7 synthesis table for end-to-end interpret
latency and kernel ratios for structural wins.  Verify suite at synthesis:
$627$ passed (series cap/ring, parse cache, lazy plots, TA incremental, order
fills, evaluator, compiler residual kernels, corpus residuals).

\subsection{Memory}

Table~\ref{tab:memory} summarizes series list memory.  Figure~\ref{fig:memory}
plots the asymptotic shape.

\begin{table}[t]
\centering
\caption{Series list memory at $n=20\,000$ bars, $s=6$ OHLCV-style lists,
$K=256$, slack $\Delta=64$.}
\label{tab:memory}
\begin{tabular}{lrrr}
\toprule
Configuration & Len/list & $\sim$pointers (total) & Ratio vs uncapped \\
\midrule
Uncapped lists & $20\,000$ & $\sim 960$k & $1\times$ \\
\code{PYNE\_SERIES\_CAP} (default) & $\le 320$ & $\sim 14.6$k & $\mathbf{\sim 78\times}$ \\
Ring $\mathrm{maxlen}=K$ & $256$ & $\sim 12.3$k & $\sim 78\times$ \\
\bottomrule
\end{tabular}
\end{table}

\subsection{Kernel speedups}

Table~\ref{tab:kernels} and Figure~\ref{fig:speedups} report incremental vs.\
full-recompute micro-benchmarks.

\begin{table}[t]
\centering
\caption{Kernel micro-benchmarks (bar-walk growing prefix, $n=2\,000$, median).}
\label{tab:kernels}
\begin{tabular}{lrrr}
\toprule
Kernel & Full recompute & Incremental & Speedup \\
\midrule
\code{kama(10)} & $1099$\,ms & $9.1$\,ms & $\mathbf{\sim 121\times}$ \\
\code{bb(20)} (pure full vs inc) & $2458$\,ms & $11.3$\,ms & $\mathbf{\sim 217\times}$ \\
\code{stochrsi(14,14)} & $10\,196$\,ms & $14.2$\,ms & $\mathbf{\sim 717\times}$ \\
\code{cmo(14)} & $13.5$\,ms & $11.4$\,ms & $\sim 1.2\times$ \\
\code{vwap} (earlier round) & $1669$\,ms & $7.7$\,ms & $\sim 217\times$ \\
\code{kc(20,2)} (earlier round) & $2054$\,ms & $26$\,ms & $\sim 78\times$ \\
\bottomrule
\end{tabular}
\end{table}

\subsection{End-to-end interpret latency}

Round~7 net benchmarks at $n=2\,000$ bars (synthesis medians):

\begin{table}[t]
\centering
\caption{Interpret Runtime median latency @ $2\,000$ bars (Round~7 synthesis).}
\label{tab:latency}
\begin{tabular}{lrrr}
\toprule
Script & R6 med (ms) & R7 med (ms) & Notes \\
\midrule
minimal & $14.92$ & $\mathbf{15.21}$ & host + one plot floor \\
\texttt{ta\_sma} & $22.53$ & $\mathbf{23.62}$ & single SMA path \\
\texttt{ta\_combo} & $137.23$ & $\mathbf{152.85}$ & plot-dominated residual \\
\texttt{strategy\_ish} & $63.07$ & $\mathbf{68.90}$ & broker + plots \\
\bottomrule
\end{tabular}
\end{table}

Kernel-level T2 wins (KAMA/BB/StochRSI) do not appear in \texttt{ta\_combo}
because that script exercises sma/ema/rsi/atr/stdev/bb/highest/lowest---already
incremental before Round~7.  Wall-time flatness of the synthesis suite is
expected; structural wins show up in kernel micros and memory, while plot
packing dominates multi-plot residual (Figure~\ref{fig:pie}).

\subsection{Compile path comparison}

Warm compile of \texttt{ta\_combo} is $\sim 0.01$\,ms with IR/disk cache; run
median $\sim 1.06$\,ms at $5\,000$ bars in nopython mode---orders of magnitude
below interpret.  Cold JIT remains highly variable (hundreds of ms to seconds)
depending on Numba cache state~\citep{lam2015numba}.  Product guidance: prefer
warm compile for multi-run TA; use interpret for compilation-ineligible scripts
and fidelity debugging.

\subsection{Feature flags summary}

\begin{center}
\begin{tabular}{lll}
\toprule
Flag & Default & Role \\
\midrule
\code{PYNE\_TA\_INCREMENTAL} & ON & interpret incremental kernels \\
\code{PYNE\_SERIES\_CAP} & ON & trim \texttt{current\_series} lists \\
\code{PYNE\_SERIES\_MAX} & unset ($K_0{=}256$) & absolute cap override \\
\code{PYNE\_SERIES\_RING} & OFF & chronological ring wrappers \\
\code{PYNE\_LIGHT\_PLOTS} & OFF & skip columnar plot capture \\
\bottomrule
\end{tabular}
\end{center}

%==============================================================================
\section{Related Work}
\label{sec:related}
%==============================================================================

\paragraph{Incremental computation.}
Classical incremental computation maintains derived results under input
changes~\citep{baker2012incremental,ramalingam1993bounded,repps1983incremental}.
Our setting is the special case of \emph{append-only time}: each bar extends
every series by one sample, so self-adjusting computation and dynamic
dependency graphs are unnecessary---fixed recurrences and sliding windows
suffice.  Streaming algorithms for sliding aggregates (sums, extrema, quantiles)
are textbook~\citep{box1970time,knuth1997taocp2}.

\paragraph{Technical analysis libraries.}
Wilder's original systems~\citep{wilder1978new}, Bollinger Bands~\citep{bollinger2001},
and Kaufman's adaptive methods~\citep{kama1998kaufman} define the mathematical
objects.  Vectorized libraries (pandas, TA-Lib, NumPy-first stacks) compute
indicators over whole arrays~\citep{harris2020numpy,brooks2015python}, which
matches \pyne{}'s compile path but not the bar-loop interpret path where
scripts interleave TA with control flow, strategy orders, and mutable
\code{var} state.

\paragraph{Memory-bounded interpreters and ring buffers.}
Ring buffers are standard in signal processing and market-data feeds.  Our
contribution is integrating them with Pine's $x[n]$ lookback semantics, $\na$
discipline, dual newest-first/chrono representations, and a flag-gated
migration that preserves corpus goldens.

\paragraph{Compiler/JIT for DSLs.}
Numba~\citep{lam2015numba} and related JITs provide the compile backend.
Partial evaluation and online compilation literature~\citep{jones1993partial}
motivates caching parse/AST/IR across multi-run product workloads.

\paragraph{Numerical accuracy.}
Higham's treatment of floating-point stability~\citep{higham2002accuracy} and
IEEE-754~\citep{ieee754} underwrite our parity budgets and sliding-variance
cautions.

%==============================================================================
\section{Limitations}
\label{sec:limits}
%==============================================================================

\begin{enumerate}[leftmargin=*]
  \item \textbf{Ring buffer default-off.}  \code{PYNE\_SERIES\_RING} remains
        opt-in until dual-write is removed and the full corpus is green under
        ring mode.
  \item \textbf{Plot-path residual.}  Incremental TA shifts the bottleneck to
        plot packing; structural wins there are only partially landed
        (\code{PYNE\_LIGHT\_PLOTS}, registry upsert hygiene).
  \item \textbf{Nested full-history leftovers.}  Some builtins (e.g., certain
        \code{obv}/\code{mode}/\code{range} paths) still recompute more than
        necessary.
  \item \textbf{Oracle quirks.}  StochRSI incremental matches the simple-RSI
        full path, not necessarily Wilder-based live platform StochRSI.
        ATR Wilder re-baseline is intentionally gated on dedicated goldens.
  \item \textbf{Full recompute + cap.}  With incremental TA disabled, recursive
        kernels can diverge from uncapped trajectories when $n\gg K$.
  \item \textbf{Host absolute variance.}  End-to-end ms figures vary by machine
        load and CPython version; treat ratios and asymptotic claims as primary.
  \item \textbf{No whole-script auto-vectorization on interpret.}  Control flow
        and strategy side effects keep the bar loop; users wanting array
        throughput should use the compile path.
\end{enumerate}

%==============================================================================
\section{Conclusion}
\label{sec:concl}
%==============================================================================

Bar-loop series languages induce characteristic memory and compute pathologies:
unbounded history and full-window TA rebuilds.  \pyne{} demonstrates that a
modest set of systems techniques---bounded series caps, chronological rings,
call-site incremental kernels, parse caching, and light plot modes---restores
linear-time, memory-bounded evaluation for the hot surface of technical
analysis while preserving Pine's lookback and $\na$ semantics.

Empirically, series caps deliver $\sim 78\times$ long-run list-memory reduction;
incremental kernels reach two--three orders of magnitude on former quadratic
rebuilds (KAMA, BB, StochRSI); and residual multi-plot interpret time is
dominated by plot packing rather than arithmetic.  The compile path, always
incremental for hot TA, remains the throughput vehicle for multi-run product
workloads.

Future work includes default-on rings after corpus ratification, deeper plot
pipeline restructuring, remaining nested-kernel residuals, and tighter dual-host
parity for worker deployments.

%==============================================================================
\section*{Acknowledgments}
%==============================================================================

We thank the \pyne{} open-source contributors and performance-agent rounds for
measurement infrastructure, goldens, and iterative kernel landings.

%==============================================================================
\appendix
\section{Memory bound (informal)}
\label{app:mem}
%==============================================================================

\begin{proposition}[Bounded series store]
\label{prop:bound}
Under \code{PYNE\_SERIES\_CAP} with capacity $K$ and slack $\Delta$, and with
$s$ chronological lists, the number of retained samples satisfies
\begin{equation}
  \forall\, t,\quad
  \sum_{j=1}^{s} |L_j(t)| \;\le\; s\,(K+\Delta).
\end{equation}
With a modular ring of maxlen $K$ per series, the bound tightens to $s\cdot K$
and each lookback of depth $<K$ is $O(1)$.
\end{proposition}

\begin{proposition}[Incremental TA time]
\label{prop:time-bound}
If every TA call site $k$ admits an update of cost $u_k\in O(1)\cup O(p_k)$
amortized, then total TA work over $n$ bars is
$O\big(n\sum_k u_k\big)$, independent of full prefix length.
\end{proposition}

\section{Environment reproduction notes}
\label{app:repro}

\begin{lstlisting}[style=python,caption={Flag surface (interpret path)}]
# Default-on incremental TA and series cap
# PYNE_TA_INCREMENTAL=1
# PYNE_SERIES_CAP=1
# Optional:
# PYNE_SERIES_MAX=512
# PYNE_SERIES_RING=1
# PYNE_LIGHT_PLOTS=1
\end{lstlisting}

\begin{lstlisting}[style=pine,caption={Minimal Pine series / TA example}]
//@version=5
indicator("Demo", overlay=true)
len = input.int(14, "Length")
v = ta.sma(close, len)
plot(v)
\end{lstlisting}

%==============================================================================
\bibliographystyle{plainnat}
\bibliography{../common/bibliography}

\end{document}
